% Generated draft: requires mathematical and source review.
\documentclass[11pt]{article}
\usepackage[margin=25mm]{geometry}
\usepackage{fontspec}
\setmainfont{Latin Modern Roman}
\usepackage{amsmath,amssymb,mathtools,unicode-math}
\setmathfont{Latin Modern Math}
\usepackage{enumitem,needspace,xurl,hyperref}
\setlength{\parindent}{0pt}
\setlength{\parskip}{5pt}
\setlength{\emergencystretch}{3em}
\newcommand{\N}{\mathbb{N}}
\newcommand{\Z}{\mathbb{Z}}
\newcommand{\Q}{\mathbb{Q}}
\newcommand{\R}{\mathbb{R}}
\newcommand{\C}{\mathbb{C}}
\newcommand{\F}{\mathbb{F}}
\newcommand{\A}{\mathbb{A}}
\newcommand{\PP}{\mathbb{P}}
\newcommand{\E}{\mathbb{E}}
\newcommand{\eps}{\varepsilon}
\newcommand{\dd}{\,\mathrm{d}}
\begin{document}
\section*{UnsolvedMath ID: 700933}
\textbf{Source problem code:} catalogue entry
\subsection{Definitions and mathematical statement}
For a set \(X\) with \(|X|\geq 2\), let \(F(X)\) be the free group on \(X\). Define its lower central series by \(\gamma_1(F(X))=F(X)\) and \(\gamma_{k+1}(F(X))=[\gamma_k(F(X)),F(X)]\). The free nilpotent group of class \(3\) on \(X\) is
\[
N_X=F(X)/\gamma_4(F(X)).
\]
The source does not specify the rank \(|X|\), whether it is finite, or the intended quantifier over ranks; this qualification remains unresolved.

A group \(G\) is torsion-free if, for every \(g\in G\) and every integer \(n\geq 1\), the equality \(g^n=1\) implies \(g=1\). A group quasiidentity is a universally quantified sentence
\[
\forall x_1\cdots\forall x_n\;\left(\bigwedge_{j=1}^{m}u_j(x_1,\ldots,x_n)=1\ \Longrightarrow\ v(x_1,\ldots,x_n)=1\right),
\]
where \(n\geq 1\), \(m\geq 0\), and \(u_1,\ldots,u_m,v\) are group words; the empty conjunction is interpreted as true. Write \(\operatorname{QI}(N_X)\) for all quasiidentities satisfied by \(N_X\).

Interpreting a basis relative to torsion-free groups in the semantic sense, a set \(\Sigma\subseteq\operatorname{QI}(N_X)\) is such a basis if, for every torsion-free group \(G\),
\[
G\models\Sigma\quad\Longleftrightarrow\quad G\models\operatorname{QI}(N_X).
\]
It is independent if, for every \(\sigma\in\Sigma\), there exists a torsion-free group \(G_\sigma\) such that
\[
G_\sigma\models\Sigma\setminus\{\sigma\},
\qquad G_\sigma\not\models\sigma.
\]
The question is whether \(N_X\) possesses an independent basis of quasiidentities relative to torsion-free groups. No finiteness restriction on the basis is stated. The displayed formulation is parameterized by \(X\); it does not resolve the source's missing rank quantifier.
\subsection{Short English statement}
Does a free nonabelian nilpotent group of class \(3\) admit a set of quasiidentities that determines all its quasiidentities among torsion-free groups, with every member indispensable within that class? The source leaves the rank and its quantification unspecified.
\subsection{Sources}
\begin{itemize}
\item Local UnsolvedMath registry: permanent ID \(700933\), source problem code \texttt{catalogue entry}, title \emph{Kourovka Notebook Problem 13.15}; supplied upstream revision \texttt{d371ee917ac9647fbb06cd6fc529deec6f3dda82}.
\item Supplied attribution: A. I. Budkin, proposed in 1995; \emph{The Kourovka Notebook (Unsolved Problems in Group Theory)}, 13th issue (1995), Problem 13.15. The supplied background also identifies the 21st issue edition (2026), arXiv:1401.0300v48. These bibliographic details are supplied source information and have not been independently verified.
\item The registry labels the problem \texttt{open}. This is a supplied status label, not an independently verified mathematical or literature assessment.
\end{itemize}
\end{document}
